\begin{textsample}{POS Dim 2 – global\_south\_2024\_11 – Score 24.00 – global\_south\_2024\_11/118190205.txt}  \label{ex:f2_pos_249}
Attachments Palestinian Territory - Despite the advent of a harsh \textbf{winter} and dire \textbf{humanitarian} circumstances, Israel continues to prevent \textbf{blankets}, clothing, and shoes -- including necessities for \textbf{children} -- from entering the Gaza Strip.
Israel has been blocking the entry of these items into the \textbf{besieged} enclave for over a year now.
As the second \textbf{winter} of Israel ' s genocidal war on the Gaza Strip begins, Palestinians are suffering from a \textbf{severe} lack of clothing and shoes, which have been banned from entering the Strip since the start of the genocide.
The only exceptions are a small number of \textbf{supplies} that are allowed in as part of \textbf{humanitarian} aid and are given to a small percentage of the roughly two million displaced people in the enclave.
Euro-Med Monitor notes that Israel restricts the entry of such items as part of its efforts to impose harsh living conditions on the Palestinian people that will ultimately lead to their actual destruction, as part of the comprehensive crime of genocide it is committing in the Gaza Strip.
There @ @ @ @ @ @ @ @ @ @ permits the prevention of basic necessities from reaching a \textbf{civilian} \textbf{population}.
Israel has destroyed at least 70 \% of the homes in the Strip and the majority of shops and markets there, including those selling clothing, in addition to limiting Palestinian merchants ' ability to coordinate the entry of goods with Israeli authorities.
Consequently, the total number of \textbf{trucks} entering the Gaza Strip in the past period contained aid that did not exceed 6 \% of the \textbf{population} ' s daily needs -- the majority of which are related to \textbf{food} \textbf{supplies} -- and the clothing and shoes allowed to enter the enclave did not exceed 0.001 \% of residents ' needs.
The vast majority of displaced people in the Gaza Strip continue to live in \textbf{tents} that do not provide adequate protection from the \textbf{cold} and rain, while hundreds of thousands of Palestinians, including women, \textbf{children}, and the elderly, are left without enough appropriate clothing to protect them from the harsh \textbf{weather} as \textbf{winter} approaches.
The lack of access to essential \textbf{medical} care in these dire circumstances also puts @ @ @ @ @ @ @ @ @ @ infections and other cold-related conditions.
The situation is made worse by the acute lack of basic medications required to treat cold-related illnesses, which is directly related to Israel ' s arbitrary blockade.
Additionally, the \textbf{population} ' s immune systems have been weakened by the scarcity of \textbf{food} and lack of variety, as well as their heavy reliance on canned \textbf{foods}, leaving them much more vulnerable than usual to viruses and illnesses.
Out of the roughly 2.3 million Palestinians living in the Gaza Strip, about two million have been forcibly displaced from their homes ; the majority of them are now living in \textbf{tents}, schools-turned-shelters, or the remains of their destroyed homes.
Those who \textbf{fled} their homes were typically forced to leave their personal belongings and clothing behind, taking only what they were wearing as they left.
Most displaced families have lost the majority of their belongings as a result of Israeli \textbf{bombardment}, and have had to search for clothing and shoes in marketplaces that have also been bombed by the occupation army.
The @ @ @ @ @ @ @ @ @ @ Strip walking barefoot in sewage- and debris-filled streets in the rain while wearing only light, shabby clothing. \textbf{Children} who lack shoes are more likely to sustain wounds and \textbf{injuries}, leaving them susceptible to infection in an environment devoid of \textbf{medical} \textbf{supplies} and medications because of the strict blockade.
People turn to short-term, unsafe, and insufficient solutions that \textbf{worsen} their suffering, like making wooden and plastic shoes for their kids.
Due to a lack of clothing, \textbf{Gazans} are currently compelled to sew or patch old clothing from old \textbf{blankets}, as only those with the means to do so can purchase any alternatives.
Due to the rainy \textbf{weather} over the past two days, the majority of the displaced have been unable to cover their \textbf{tents} and protect them from the rain, which has resulted in hundreds of \textbf{tents} flooding and the few belongings of the displaced becoming drenched in \textbf{water}.
Notably, Israel also prohibits the entry of adequate quantities of \textbf{tents}, tarps, and nylon into the Strip, as well as other necessities @ @ @ @ @ @ @ @ @ @, firewood, \textbf{fuel}, and heating sources.
Israel ' s continuous and \textbf{severe} deprivation of the fundamental necessities of life is an act of genocide, as it seeks to strip the Palestinian \textbf{population} of the most basic means of protection, with the aim of physically erasing their existence. \textbf{Children} and other vulnerable groups are specifically targeted by Israel as they are more affected by this deprivation, which exacerbates their suffering and raises the death \textbf{rates} among them ; due to the lack of \textbf{refuge} from \textbf{winter} \textbf{weather}, these \textbf{rates} will undoubtedly spike without international intervention.
Denying basic necessities to all segments of the \textbf{civilian} \textbf{population} is an outright assault on people ' s dignity and a deprivation of their humanity.
Treating them as though they are undeserving of even the most basic rights has shattered their spirits, contributing to a sense of dejection felt by all Gaza Strip residents.
In creating such inhuman conditions, Israel also expresses a clear aim to destroy Palestinians ' cultural and social identity.
International and United Nations organisations must work, @ @ @ @ @ @ @ @ @ @ the entry of basic materials into the Gaza Strip, and to publicly expose these crimes.
Given the grave worsening of the \textbf{humanitarian} situation, the international community must take responsibility for halting the genocide in the Gaza Strip and all related crimes being committed by Israel and its allies, as this is the only way to protect \textbf{civilians} and preserve what remains.
In addition to imposing sanctions on Israel and implementing the arrest warrants issued by the International Criminal Court against the Israeli Prime Minister and Minister of Defense as soon as possible, as well as their transfer to international custody, it is imperative that Palestinians in the Gaza Strip be given immediate and unhindered access to \textbf{winter} clothing, shoes, and the most basic tools of survival.

% matched lemmas: besieged, blanket, bombardment, child, civilian, cold, flee, food, fuel, gazan, humanitarian, injury, medical, population, rate, refuge, severe, supply, tent, truck, water, weather, winter, worsen
\end{textsample}
